% GENERATED FILE. Edit canonical lessons, then run npm run generate:books.
% canonical-source-hash: fnv1a64:7708c811d71512d2
% canonical-lessons: GE-C415-schliessfach, GE-C415-gepaeckaufbewahrung, GE-C415-rolltreppe, GE-C415-schwarzfahren, GE-C415-anstellen

\chapter{Am großen Bahnhof}
\label{ch:ge-am-gro-en-bahnhof}
\hlchaptermodality{415}

\begin{chapteropening}
\textbf{By the end of this chapter:} \emph{I can manage at a busy station: the locker, the left-luggage office, the escalator, riding without a ticket and queueing.}

Complete the last lesson of chapter 415: I can manage at a busy station: the locker, the left-luggage office, the escalator, riding without a ticket and queueing.
\end{chapteropening}

\section[das Schließfach]{das Schließfach — the locker}

\label{lesson:GE-C415-schliessfach}



\begin{quote}
Before the new one: say the German for to commute, then the German for the commuter.
\end{quote}

\subsection*{You'll want to know: das Schließfach}
\textbf{das Schließfach} — \textquotedblleft{}the locker\textquotedblright{}.

\begin{quote}
Er legt seinen Rucksack für zwei Stunden in ein Schließfach.
\end{quote}

\begin{quote}
\emph{He puts his rucksack in a locker for two hours.}
\end{quote}

Plural: \emph{die Schließfächer}. Also a safe-deposit box in a bank.

\begin{cousinweb}
\emph{Schließfach} is built from \emph{schließen} 'to lock' and \emph{Fach} 'compartment', both inherited: \emph{Fach} goes back to Old High German \emph{fah} (EtymDB).
\end{cousinweb}

\begin{grammarlens}[title={What you've built}]
One more word for telling, explaining and describing at B1.
\end{grammarlens}

\subsection*{Your turn}
\begin{itemize}
\raggedright
  \item \emph{Say it:} \emph{das Schließfach}
  \item \emph{Say it:} \emph{das Schließfach}, once more
  \item \emph{Say it:} a sentence of your own with \emph{das Schließfach}
  \item \emph{Recall:} say the German for to commute, then the German for the commuter, then say \emph{das Schließfach} again
\end{itemize}

\begin{tcolorbox}[breakable,colback=teal!4,colframe=teal!35!black,title={Before you move on}]
What does das Schließfach mean? (\textquotedblleft{}The locker\textquotedblright{}.) Say it once more.
\end{tcolorbox}
\section[die Gepäckaufbewahrung]{die Gepäckaufbewahrung — the left-luggage office}

\label{lesson:GE-C415-gepaeckaufbewahrung}



\begin{quote}
Before the new one: say the German for the commuter, then the German for the locker.
\end{quote}

\subsection*{You'll want to know: die Gepäckaufbewahrung}
\textbf{die Gepäckaufbewahrung} — \textquotedblleft{}the left-luggage office\textquotedblright{}.

\begin{quote}
Die Gepäckaufbewahrung ist neben dem Ausgang.
\end{quote}

\begin{quote}
\emph{The left-luggage office is next to the exit.}
\end{quote}

No plural in everyday use. \emph{Aufbewahren} — to keep, to store.

\begin{cousinweb}
\emph{Gepäckaufbewahrung} is built from \emph{Gepäck} 'luggage', from \emph{packen} 'to pack', and \emph{Aufbewahrung} 'keeping', from \emph{aufbewahren} 'to keep safe', which is built on \emph{wahren} 'to guard'.
\end{cousinweb}

\begin{grammarlens}[title={What you've built}]
One more word for telling, explaining and describing at B1.
\end{grammarlens}

\subsection*{Your turn}
\begin{itemize}
\raggedright
  \item \emph{Say it:} \emph{die Gepäckaufbewahrung}
  \item \emph{Say it:} \emph{die Gepäckaufbewahrung}, once more
  \item \emph{Say it:} a sentence of your own with \emph{die Gepäckaufbewahrung}
  \item \emph{Recall:} say the German for the commuter, then the German for the locker, then say \emph{die Gepäckaufbewahrung} again
\end{itemize}

\begin{tcolorbox}[breakable,colback=teal!4,colframe=teal!35!black,title={Before you move on}]
What does die Gepäckaufbewahrung mean? (\textquotedblleft{}The left-luggage office\textquotedblright{}.) Say it once more.
\end{tcolorbox}
\section[die Rolltreppe]{die Rolltreppe — the escalator}

\label{lesson:GE-C415-rolltreppe}



\begin{quote}
Before the new one: say the German for the locker, then the German for the left-luggage office.
\end{quote}

\subsection*{You'll want to know: die Rolltreppe}
\textbf{die Rolltreppe} — \textquotedblleft{}the escalator\textquotedblright{}.

\begin{quote}
Die Rolltreppe zum Bahnsteig ist kaputt.
\end{quote}

\begin{quote}
\emph{The escalator to the platform is broken.}
\end{quote}

Plural: \emph{die Rolltreppen}.

\begin{cousinweb}
\emph{Rolltreppe} is built from \emph{rollen} 'to roll', which goes back through French to Latin \emph{rotula} 'little wheel', and \emph{Treppe} 'stairs', an inherited word that came into the standard language from Low German (EtymDB lists \emph{Rolltreppe} under \emph{Treppe}).
\end{cousinweb}

\begin{grammarlens}[title={What you've built}]
One more word for telling, explaining and describing at B1.
\end{grammarlens}

\subsection*{Your turn}
\begin{itemize}
\raggedright
  \item \emph{Say it:} \emph{die Rolltreppe}
  \item \emph{Say it:} \emph{die Rolltreppe}, once more
  \item \emph{Say it:} a sentence of your own with \emph{die Rolltreppe}
  \item \emph{Recall:} say the German for the locker, then the German for the left-luggage office, then say \emph{die Rolltreppe} again
\end{itemize}

\begin{tcolorbox}[breakable,colback=teal!4,colframe=teal!35!black,title={Before you move on}]
What does die Rolltreppe mean? (\textquotedblleft{}The escalator\textquotedblright{}.) Say it once more.
\end{tcolorbox}
\section[schwarzfahren]{schwarzfahren — to ride without a ticket, to dodge the fare}

\label{lesson:GE-C415-schwarzfahren}



\begin{quote}
Before the new one: say the German for the left-luggage office, then the German for the escalator.
\end{quote}

\subsection*{You'll want to know: schwarzfahren}
\textbf{schwarzfahren} — \textquotedblleft{}to ride without a ticket, to dodge the fare\textquotedblright{}.

\begin{quote}
Ein junger Mann fährt schwarz und muss sechzig Euro zahlen.
\end{quote}

\begin{quote}
\emph{A young man rides without a ticket and has to pay sixty euros.}
\end{quote}

Separable; Perfekt with \emph{sein}: \emph{er ist schwarzgefahren}. \emph{Der Schwarzfahrer} — the fare dodger.

\begin{cousinweb}
\emph{schwarzfahren} is built from \emph{schwarz} 'black', in its old sense 'secret, against the rules', and the inherited verb \emph{fahren} (EtymDB). Dutch \emph{zwartrijden} is built the same way.
\end{cousinweb}

\begin{grammarlens}[title={What you've built}]
One more word for telling, explaining and describing at B1.
\end{grammarlens}

\subsection*{Your turn}
\begin{itemize}
\raggedright
  \item \emph{Say it:} \emph{schwarzfahren}
  \item \emph{Say it:} \emph{schwarzfahren}, once more
  \item \emph{Say it:} a sentence of your own with \emph{schwarzfahren}
  \item \emph{Recall:} say the German for the left-luggage office, then the German for the escalator, then say \emph{schwarzfahren} again
\end{itemize}

\begin{tcolorbox}[breakable,colback=teal!4,colframe=teal!35!black,title={Before you move on}]
What does schwarzfahren mean? (\textquotedblleft{}To ride without a ticket, to dodge the fare\textquotedblright{}.) Say it once more.
\end{tcolorbox}
\section[anstellen]{anstellen — to queue (mostly reflexive: sich anstellen); to switch on}

\label{lesson:GE-C415-anstellen}



\begin{quote}
Before the new one: say the German for the escalator, then the German for to ride without a ticket.
\end{quote}

\subsection*{You'll want to know: anstellen}
\textbf{anstellen} — \textquotedblleft{}to queue (mostly reflexive: sich anstellen); to switch on\textquotedblright{}.

\begin{quote}
Vor dem Fahrkartenautomaten stellen wir uns hinten an.
\end{quote}

\begin{quote}
\emph{At the ticket machine we join the back of the queue.}
\end{quote}

Mostly reflexive and separable: \emph{ich stelle mich an}. \emph{Das Radio anstellen} — to switch the radio on.

\begin{cousinweb}
\emph{anstellen} is built from \emph{an} 'at, on' and \emph{stellen} 'to place, to set', inherited from Old High German \emph{stellen} (EtymDB): to place oneself at the end of a line.
\end{cousinweb}

\begin{grammarlens}[title={What you've built}]
One more word for telling, explaining and describing at B1.
\end{grammarlens}

\subsection*{Your turn}
\begin{itemize}
\raggedright
  \item \emph{Say it:} \emph{anstellen}
  \item \emph{Say it:} \emph{anstellen}, once more
  \item \emph{Say it:} a sentence of your own with \emph{anstellen}
  \item \emph{Recall:} say the German for the escalator, then the German for to ride without a ticket, then say \emph{anstellen} again
\end{itemize}

\begin{tcolorbox}[breakable,colback=teal!4,colframe=teal!35!black,title={Before you move on}]
What does anstellen mean? (\textquotedblleft{}To queue (mostly reflexive: sich anstellen); to switch on\textquotedblright{}.) Say it once more.
\end{tcolorbox}
